\chapter*{Tóm tắt}
\addcontentsline{toc}{chapter}{Tóm tắt}
Báo cáo này trình bày quy trình nghiên cứu, xây dựng và đánh giá thực nghiệm hệ thống hỏi đáp trừu tượng tiếng Việt trong miền y khoa dựa trên bộ dữ liệu ViMedAQA. Khác với nhiệm vụ hỏi đáp trích xuất truyền thống, hỏi đáp trừu tượng y khoa đòi hỏi mô hình phải hiểu sâu sắc ngữ cảnh và sinh câu trả lời tự nhiên, trôi chảy và chính xác về mặt y khoa. Chúng tôi thực hiện nghiên cứu so sánh hai kiến trúc sinh chuỗi chuyên biệt cho tiếng Việt là ViT5-base và BARTpho-word bằng cách tinh chỉnh (fine-tuning) trên tập dữ liệu gồm 23.553 mẫu. Đồng thời, mô hình ngôn ngữ lớn Llama-3.3-70B-versatile được tích hợp làm đường cơ sở (baseline) dưới dạng zero-shot thông qua Groq API để làm hệ quy chiếu đánh giá. Kết quả kiểm thử định lượng cho thấy cả hai mô hình tinh chỉnh đều đạt hiệu năng vượt trội. Cụ thể, ViT5-base đạt điểm ROUGE-L là 66,80\%, BLEU-4 là 50,05\% và BERTScore F1 là 87,52\%, trong khi BARTpho-word đạt ROUGE-L là 66,81\% và BERTScore F1 là 82,30\%. Đáng chú ý, baseline Llama-3.3 zero-shot thể hiện khả năng vượt trội với BERTScore F1 đạt 87,57\% nhờ tri thức nền tảng khổng lồ. Báo cáo cũng phân tích sâu sắc các nhóm lỗi sinh văn bản thực tế (lặp từ, ảo tưởng thông tin y khoa, độ lệch chiều dài) và nghiên cứu sự ảnh hưởng của chiến lược giải mã (beam search giúp cải thiện 0,6\% điểm ROUGE so với greedy decoding). Cuối cùng, một ứng dụng minh họa Gradio tích hợp cơ chế truy xuất thông tin BM25 (RAG) đã được triển khai và công bố công khai. Kết quả nghiên cứu đóng góp một pipeline hoàn chỉnh, có khả năng tái lập cao cho các ứng dụng y tế thông minh bằng tiếng Việt.
\clearpage